% ===================================================================
%  TAULA N° 10 — La Mar. — Lo Pòrt.
%  Traduction 2026 d'apres texto/io, controlee sur texto/fr et
%  texto/ca. Consigne : voir texto/oc/10-taula-01.tex.
%
%  Le francais decale sa numerotation d'un cran a partir de l'alinea
%  8 ; c'est une coquille de l'atelier francais, et la colonne suit
%  l'ido, qui compte juste.
% ===================================================================

%%K t10-apar-1 apar
\VUsaut{4.0mm}
\VUtitre{15.0pt}{60}{TAULA N° 10}
\VUsaut{3.0mm}
\VUpk{11.4pt}{40}{La Mar. --- Lo Pòrt.}
\VUsaut{3.0mm}

%%K t10-01-1 p
1. --- L'estiu, mentre que d'unes van cercar la calma e la frescor a la montanha,
d'autres aiman mai anar a la riba de la mar. Aquò es çò que faguèrem l'an passat,
que nos agrada fòrça l'\VUgras{ocean} \textsuperscript{(1)}.

%%K t10-02-1 p
2. --- Per ma part, senti un grand plaser a contemplar aquela immensa napa d'aiga
que s'espandís fins a l'\VUgras{orizont} \textsuperscript{(2)} e a veire partir
per una longa \VUgras{traversada} los enòrmes \VUgras{naviris de vapor}
\textsuperscript{(3)} ont embarcan los viatjaires que van en America.

%%K t10-03-1 p
3. --- Per aquò mon paire, que coneis mos gostes, causís totjorn, per passar las
vacanças, una plaja fòrça tranquilla situada contra un de nòstres grands
\VUgras{pòrts de guèrra}.

%%K t10-03-2 p
Pendent nòstra darrièra estada, l'\VUgras{esquadra} èra venguda getar l'ancora
darrièr l'enòrme \VUgras{dic} \textsuperscript{(4)} que barra e protegís lo
\VUgras{pòrt militar} \textsuperscript{(5)}, e poguèri admirar a mon gost los
diferents \VUgras{naviris de guèrra}, del pichon \VUgras{soslevatri}
\textsuperscript{(10)}, que se somergís e despareis jos las aigas, fins a l'enòrme
\VUgras{cuirassat} \textsuperscript{(9)}, que fins ara cranhiá gaireben pas que
los assauts del \VUgras{torpedaire} \textsuperscript{(8)}.

%%K t10-tit-2 sub
\VUsaut{3.0mm}
\VUpk{10.2pt}{40}{Las embarcacions pichonas.}
\VUsaut{3.0mm}

%%K t10-04-1 p
4. --- Aqueste sabiá ja descobrir fòrça \VUgras{abilament} lo punt feble de la
gròssa cuirassa metallica per i clavar lo perilhós \VUgras{torpedò}, que son
explosion causa la pèrda del naviri; mas èra mens de crénher que lo soslevatri,
que los contratorpedaires li donavan caça aisidament.

%%K t10-05-1 p
5. --- Poguèri tanben veire los rapids \VUgras{crosaires} \textsuperscript{(6)}
cuirassats, gaireben tan invulnerables coma lo veritable cuirassat, mantun
\VUgras{transpòrt} \textsuperscript{(12)}, tres o quatre \VUgras{canonièras}
\textsuperscript{(7)} e, enfin, una \VUgras{fragata} \textsuperscript{(11)}.
\VUgras{L'espectacle d'aquela flòta, quand manòbra per levar l'ancora, es dels mai
interessants.}

%%K t10-06-1 p
6. --- Quina diferéncia de construccion entre aquelas enòrmas massas d'acièr e los
elegants \VUgras{bricbarcas} o los prims \VUgras{velièrs}
\textsuperscript{(13)}, que s'emplegan encara dins la marina mercanda, e que ieu
vesiái dintrar al pòrt a tota vela quand anavi passejar sul \VUgras{cai}
\textsuperscript{(14)}! Es verai qu'aqueles perdián fòrça de lors avantatges quand
lo \VUgras{vent} èra contrari. Lor caliá recampar totas las velas per dintrar dins
la darsena, e caliá un \VUgras{remolcador} \textsuperscript{(16)} qu'anèsse los
quèrre per los menar al cai coma de simplas \VUgras{gabarras}
\textsuperscript{(17)}.

%%K t10-tit-3 sub
\VUsaut{3.0mm}
\VUpk{10.2pt}{40}{Lo pòrt de comèrci.}
\VUsaut{3.0mm}

%%K t10-07-1 p
7. --- Çò qu'es interessant dins lo pòrt que ne parli es qu'apara pas sonque un
pòrt de guèrra, creat artificialament amb d'obratges gigantescs, mas en mai un
meravilhós \VUgras{pòrt de comèrci} situat a l'\VUgras{embocadura} d'un grand
riu.

%%K t10-08-1 p
8. --- La lenga de tèrra que separa las doas darsenas es fortificada e defenduda
per una seguida de \VUgras{bateriás} e de \VUgras{bastions} que se fican dins la
mar. Cal una permission especiala per passejar suls \VUgras{ramparts}
\textsuperscript{(15)}. Ieu l'aviái obtenguda sens pena per mejan de mon oncle,
qu'es enginhaire dels \VUgras{cantièrs} \textsuperscript{(18)}, e anèri visitar
los naviris que se bastisson jos de vasts cobèrts.

%%K t10-09-1 p
9. --- Assistiguèri quitament al lançament d'un cuirassat e me sovèni del moment
d'emocion que sentiguèt tota la multitud quand l'enòrma massa, abandonada a
ela meteissa, comencèt de lisar sus sa cala e venguèt susnadar sus l'aiga del riu.

%%K t10-10-1 p
10. --- Per anar als cantièrs se crosa lo \VUgras{camp de manòbras}
\textsuperscript{(19)}, ont las tropas de la garnison venon cada jorn far
l'exercici, e se passa per una \VUgras{passarèla} \textsuperscript{(20)} de fèrre
qu'unís l'esplanada amb l'\VUgras{arsenal} \textsuperscript{(21)}.

%%K t10-tit-4 sub
\VUsaut{3.0mm}
\VUpk{10.2pt}{40}{L'Arsenal e los dòcs.}
\VUsaut{3.0mm}

%%K t10-11-1 p
11. --- Amb mon oncle visitèri aquel grand establiment plen de \VUgras{canons}
\textsuperscript{(22)}, de \VUgras{balas} \textsuperscript{(23)} e
d'\VUgras{obuses}. Vegèri tanben la \VUgras{fabrica metallurgica}
\textsuperscript{(24)} ont se fabrican las \VUgras{caudièras}
\textsuperscript{(25)} dels naviris de vapor.

%%K t10-12-1 p
12. --- Fòrça prèp i a los \VUgras{dòcs} \textsuperscript{(26)} --- dòcs de carena
--- e los curiosissims \VUgras{dòcs flotants} \textsuperscript{(27)}, ont poguèri
veire de prèp, sus un naviri en reparacion, l'\VUgras{elici}
\textsuperscript{(28)}, la \VUgras{quilha} \textsuperscript{(29)} e lo
\VUgras{timon} \textsuperscript{(30)} d'una nau.

%%K t10-13-1 p
13. --- Lo pòrt de comèrci es tan interessant d'estudiar coma lo pòrt de guèrra,
e passèri fòrça oras a badar suls cais ont son amarrats los \VUgras{vapors de
ròdas} \textsuperscript{(31)} que remontan lo riu, e a agachar foncionar la
\VUgras{grua a mastrejar} \textsuperscript{(32)}, qu'auça coma de plumas de cargas
enòrmas.

%%K t10-tit-5 sub
\VUsaut{4.0mm}
\VUpk{10.2pt}{40}{Lo Pilòt.}
\VUsaut{3.0mm}

%%K t10-14-1 p
14. --- Admiravi los \VUgras{transatlantics} \textsuperscript{(34)} que sortissián
de la rada amb lors \VUgras{bandièras} \textsuperscript{(35)} al vent, menats per
un abil \VUgras{pilòt} \textsuperscript{(36)} que sabiá meravilhosament seguir lo
\VUgras{canal} marcat per de \VUgras{boiàs} \textsuperscript{(40)} e de
\VUgras{balisas}, en evitar las \VUgras{gabarras} \textsuperscript{(41)} que se
menan amb tanta pena amb lor sola vela carrada. Destriavi a la \VUgras{prova}
\textsuperscript{(37)} las \VUgras{cabinas} \textsuperscript{(39)} de segonda
classa e, a la \VUgras{popa} \textsuperscript{(38)}, los salons dels passatgièrs
de primièra.

%%K t10-15-1 p
15. --- De còps assistissiái tanben al cargament d'un \VUgras{naviri de carga}
\textsuperscript{(42)}, ont s'amolonavan las mercadarias del \VUgras{magazin}
\textsuperscript{(43)} e de l'\VUgras{estacion maritima} \textsuperscript{(44)},
portadas pels nombroses trens de la \VUgras{linha dels cais}
\textsuperscript{(45)}.

%%K t10-tit-6 sub
\VUsaut{3.0mm}
\VUpk{10.2pt}{40}{Ramar per plaser.}
\VUsaut{3.0mm}

%%K t10-16-1 p
16. --- Sovent ramavi dins la \VUgras{darsena a flòt} \textsuperscript{(53)}, ont
se refugian las pichonas \VUgras{embarcacions} \textsuperscript{(46)}, e èra una
jòia per ieu de manejar lo \VUgras{rem} \textsuperscript{(47)}. Pujavi sul
\VUgras{pont} \textsuperscript{(48)} d'una \VUgras{barca de vela}
\textsuperscript{(49)}, m'enfilavi amb l'ajuda del \VUgras{cordatge}
\textsuperscript{(52)} pel \VUgras{mast} \textsuperscript{(51)} fins a las
\VUgras{vergas} \textsuperscript{(50)}, e puèi reprenia mon passejament en
\VUgras{iòla} \textsuperscript{(54)} entre las pesucas \VUgras{barcas de pesca}
\textsuperscript{(55)}, ont s'eissugan los \VUgras{filats}
\textsuperscript{(56)}.

%%K t10-17-1 p
17. --- Sul \VUgras{cai} \textsuperscript{(58)} desfilavan davant ieu las
\VUgras{mercadarias} \textsuperscript{(59)} mai divèrsas, ficadas dins de
\VUgras{caissas} \textsuperscript{(60)} o dins de \VUgras{balas}
\textsuperscript{(61)}. Formavan lo \VUgras{cargament} \textsuperscript{(63)} de
las barcassas, e de poderosas \VUgras{gruas} \textsuperscript{(62)} las auçavan e
las pausavan dins de \VUgras{camions} \textsuperscript{(64)} mentre que los
\VUgras{transportaires} fasián la declaracion e pagavan los drechs a la
\VUgras{doana} \textsuperscript{(65)}.

%%K t10-18-1 p
18. --- Enfin, quand arribavi al \VUgras{pont giratòri} \textsuperscript{(66)} o a
l'\VUgras{esclusa} \textsuperscript{(69)}, en passar davant la \VUgras{draga}
\textsuperscript{(68)} que neteja lo \VUgras{canal} \textsuperscript{(67)},
desembarcavi sul cai, contra lo \VUgras{semafòr} \textsuperscript{(70)}.

%%K t10-19-1 p
19. --- Es de l'autre costat d'aquel cai que s'acòstan las \VUgras{chalupas}
\textsuperscript{(71)} que menan a tèrra los oficièrs de l'esquadra. Es un
espectacle bèl de veire los mariniers auçar lors rems a la mesura mentre que
l'embarcacion, portada per l'\VUgras{onda} \textsuperscript{(72)}, abòrda sens
esfòrç l'escala del desembarcador. Es aicí que se pòt melhor observar la
\VUgras{marèa} e lo movement del \VUgras{flux} e del \VUgras{reflux} sus la
\VUgras{plaja} \textsuperscript{(73)}.

%%K t10-tit-7 sub
\VUsaut{3.0mm}
\VUpk{10.2pt}{40}{Lo Cercle dels oficièrs.}
\VUsaut{3.0mm}

%%K t10-20-1 p
20. --- Per tornar a l'ostal prenia lo \VUgras{tramvai electric}
\textsuperscript{(74)}, que son \VUgras{arrèst} \textsuperscript{(75)} es al
\VUgras{pal} qu'es davant lo cercle dels oficièrs.

%%K t10-20-2 p
Del \VUgras{balcon} \textsuperscript{(76)} del cercle, ornat d'\VUgras{ancoras}
\textsuperscript{(77)}, de canons e de balas, vesiái sovent una brilhanta
acampada d'oficièrs de marina: de \VUgras{tenents de vaissèl}
\textsuperscript{(78)} amb lor espasa, de \VUgras{capitanis d'artilhariá}
\textsuperscript{(79)} e d'\VUgras{infantariá} \textsuperscript{(80)} de marina
amb lor \VUgras{sabre}. Darrièr eles i aviá un \VUgras{marinièr}
\textsuperscript{(81)} que lor portava una \VUgras{lunèta} quand volián destriar
çò que se passava al luènh.

%%K t10-21-1 p
21. --- Vesiái tanben de joves \VUgras{alfièrs de vaissèl}
\textsuperscript{(82)} que recebián los òrdres de lor \VUgras{comandant}
\textsuperscript{(83)} o de l'\VUgras{almirant} \textsuperscript{(84)}, mentre
qu'un \VUgras{caporal} \textsuperscript{(85)} esperava per los portar a bòrd.

%%K t10-22-1 p
22. --- D'aquel balcon la vista es magnifica: se domina tota la plaja amb sas
\VUgras{tendas} \textsuperscript{(86)} e sas \VUgras{cabanas}
\textsuperscript{(87)}, e se vei, a l'ora del banh, los \VUgras{banhaires}
\textsuperscript{(88)} que jògan alègrament davant lo \VUgras{casinò}
\textsuperscript{(89)}.

%%K t10-23-1 p
23. --- A la cima de la plaja, la còsta forma una pichona \VUgras{peninsula}
\textsuperscript{(91)} rocalhosa, ont ven se copar la \VUgras{resaca}
\textsuperscript{(92)} e sus la quala es bastit un \VUgras{far}
\textsuperscript{(93)} qu'indica de nuèch lo camin als navegants. Aquela peninsula
es unida a la tèrra pas que per un \VUgras{istme} \textsuperscript{(90)} fòrça
estrech.

%%K t10-tit-8 sub
\VUsaut{3.0mm}
\VUpk{10.2pt}{40}{Vista generala de la còsta.}
\VUsaut{3.0mm}

%%K t10-24-1 p
24. --- La \VUgras{ròca talhada} \textsuperscript{(94)} se perlonga, retalhada per
la mar, que i dessenha capriciosament una seguida de \VUgras{golfes}
\textsuperscript{(95)} e de \VUgras{promontòris} (caps) que la fan dels mai
pintorescs.

%%K t10-24-2 p
Sul \VUgras{planòl} \textsuperscript{(96)}, que domina l'\VUgras{estrech}
\textsuperscript{(98)}, i a un \VUgras{fòrt} \textsuperscript{(97)} fòrça
important e qualques vilas.

%%K t10-25-1 p
25. --- Aquel estrech separa del continent una \VUgras{illa}
\textsuperscript{(99)} fòrça perilhosa, envirotada d'\VUgras{esculhs}
\textsuperscript{(100)} e de \VUgras{rescluses}, ont de barcas e quitament de
grands naviris s'encalan de còps. Malgrat la prontor dels socors e l'arribada
rapida de las \VUgras{barcas de sauvatge}, aqueles \VUgras{naufragis} son
terribles e, dins las fòrtas \VUgras{tempèstas} d'equinòxi, mantun naviri s'es
perdut amb òmes e mercadarias.

%%K t10-25-2 p
Malgrat aquel vesinatge, \VUgras{lo pòrt} lo frequentan fòrça los mariniers.
\VUsaut{6.0mm}
\VUfilet{22mm}
